% META: {"id":"TCOMPL-AIR-CO-003","chapitre":"TCOMPL-CALCULS-AIRES","type_objet":"corrige","exercice_ref":"TCOMPL-AIR-EX-003","capacites_codes":["C3"],"sources_inspiration":[],"status":"approved","origin":"generated","status_history":["generated","audited","accepted","approved"],"release_acceptance":"RELEASE_OWNER_BATCH_ACCEPTANCE","acceptance_closure_digest":"sha256:9b3ccf9a81c5520fbb7e03b7b2d3e7bf2057a02834b6d908bf3be5f49f3b553f","acceptance_receipt_digest":"sha256:4fad86f0282e0fa4e0419cca1ad6379ae7af5a280c04a4a90880f90a1c044242"}
% BEGIN-VERIFY
% from sympy import symbols, solve, Eq, integrate
% x = symbols('x')
% f = -x**2 + 4
% g = x**2 - 2*x
% inter = solve(Eq(f, g), x)
% assert inter == [-1, 2]
% aire = integrate(f - g, (x, -1, 2))
% assert aire == 9
% END-VERIFY
\begin{corrige}{TCOMPL-AIR-EX-003}
\textbf{1.} $f(x)=g(x) \iff -x^2+4 = x^2-2x \iff -2x^2+2x+4=0 \iff x^2-x-2=0$. Discriminant $\Delta=1+8=9$, racines $x=\dfrac{1\pm3}{2}$, soit $x=2$ ou $x=-1$.

\textbf{2.} $f(x)-g(x) = -2x^2+2x+4 = -2(x+1)(x-2)$, qui est un polynôme du second degré de coefficient dominant négatif, positif \textbf{entre} ses racines : $f(x)-g(x)\geq0$ sur $[-1,2]$, donc $f\geq g$ sur cet intervalle.

\textbf{3.} L'aire vaut :
\[
  \int_{-1}^{2} \big(f(x)-g(x)\big)\,\mathrm{d}x = \int_{-1}^{2}\big(-2x^2+2x+4\big)\,\mathrm{d}x = \left[-\frac{2x^3}{3}+x^2+4x\right]_{-1}^{2} = 9 \text{ unités d'aire.}
\]
\end{corrige}
